\chapter{1977年江苏常州市初试卷}

\section{解答题}

\begin{problem}
计算：\(-1-\left(-5\frac{1}{2}\right)\times\frac{4}{11}+\left(-2\right)^3\div\left[\left(-3\right)^2+7\right]\).
\end{problem}

\begin{solution}
先把带分数、幂和括号内的数分别化简：\(-5\frac{1}{2}=-\frac{11}{2}\)，\((-2)^3=-8\)，\((-3)^2+7=9+7=16\). 因此
\[
-1-\left(-\frac{11}{2}\right)\times\frac{4}{11}+\left(-2\right)^3\div\left[\left(-3\right)^2+7\right]
=-1-(-2)+\frac{-8}{16}=1-\frac{1}{2}=\frac{1}{2}
\]
故原式的值为 \(\frac{1}{2}\).
\end{solution}

\begin{problem}
计算：\(-\frac{2}{2+\sqrt{2}}+\frac{1}{\sqrt{3}-\sqrt{2}}-\frac{3+\sqrt{3}}{3-\sqrt{3}}\).
\end{problem}

\begin{solution}
分别乘以各分母的共轭式进行有理化.

第一项为
\[
-\frac{2}{2+\sqrt{2}}=-\frac{2(2-\sqrt{2})}{(2+\sqrt{2})(2-\sqrt{2})}=-\frac{2(2-\sqrt{2})}{4-2}=\sqrt{2}-2
\]
第二项为
\[
\frac{1}{\sqrt{3}-\sqrt{2}}=\frac{\sqrt{3}+\sqrt{2}}{(\sqrt{3}-\sqrt{2})(\sqrt{3}+\sqrt{2})}=\frac{\sqrt{3}+\sqrt{2}}{3-2}=\sqrt{3}+\sqrt{2}
\]
第三项为
\[
\frac{3+\sqrt{3}}{3-\sqrt{3}}=\frac{(3+\sqrt{3})^2}{(3-\sqrt{3})(3+\sqrt{3})}=\frac{9+6\sqrt{3}+3}{9-3}=2+\sqrt{3}
\]
把三项代回原式，注意第三项前有负号，得
\[
=\left(\sqrt{2}-2\right)+\left(\sqrt{3}+\sqrt{2}\right)-\left(2+\sqrt{3}\right)=2\sqrt{2}-4
\]
故原式的值为 \(2\sqrt{2}-4\).
\end{solution}

\begin{problem}
计算：\(\log_2\left(1+\lg 10\right)\).
\end{problem}

\begin{solution}
由常用对数的定义，\(\lg 10=1\)，所以括号内的数为 \(1+\lg 10=1+1=2\). 又因为 \(2^1=2\)，根据对数的定义可得
\(\log_2\left(1+\lg 10\right)=\log_2 2=1\).
故原式的值为 \(1\).
\end{solution}

\begin{problem}
\(\sin 390^\circ\cos\left(-60^\circ\right)+\sin 120^\circ\cos 750^\circ\).
\end{problem}

\begin{solution}
利用正弦、余弦的周期性以及余弦函数的偶性，逐项化为锐角三角函数值：
\(\sin 390^\circ=\sin\left(360^\circ+30^\circ\right)=\frac{1}{2}\)，
\(\cos\left(-60^\circ\right)=\cos 60^\circ=\frac{1}{2}\)，
\(\sin 120^\circ=\sin\left(180^\circ-60^\circ\right)=\frac{\sqrt{3}}{2}\)，
\(\cos 750^\circ=\cos\left(720^\circ+30^\circ\right)=\frac{\sqrt{3}}{2}\). 因此原式为
\[
=\frac{1}{2}\cdot\frac{1}{2}+\frac{\sqrt{3}}{2}\cdot\frac{\sqrt{3}}{2}
=\frac{1}{4}+\frac{3}{4}=1
\]
故原式的值为 \(1\).
\end{solution}

\begin{problem}
化简：\(a+\sqrt{a^2-2a+1}\).
\end{problem}

\begin{solution}
对任意 \(a\in\R\)，被开方数可配方为 \(a^2-2a+1=(a-1)^2\)，而算术平方根满足 \(\sqrt{u^2}=\lvert u\rvert\). 因此 \(\sqrt{a^2-2a+1}=\sqrt{(a-1)^2}=\lvert a-1\rvert\)，原式化为 \(a+\lvert a-1\rvert\).

绝对值的符号由 \(a-1\) 的正负决定. 当 \(a\geq 1\) 时，\(\lvert a-1\rvert=a-1\)，原式为 \(a+a-1=2a-1\)；当 \(a<1\) 时，\(\lvert a-1\rvert=1-a\)，原式为 \(a+1-a=1\). 所以
\[
=\begin{cases}
2a-1,&a\geq 1,\\
1,&a<1.
\end{cases}
\]
在边界 \(a=1\) 时两种表达式均为 \(1\)，故分类完整.
\end{solution}

\begin{problem}
已知 \(a=9\)，\(b=\sqrt{3}\)，求 \(a^{\frac{3}{2}}b\left(\frac{b^2}{a}\right)^{\frac{1}{2}}\div\left(\frac{1}{a}\right)^{-\frac{1}{2}}\) 的值.
\end{problem}

\begin{solution}
代入 \(a=9\)，\(b=\sqrt{3}\). 因为 \(a>0\)，有 \(a^{\frac{3}{2}}=9^{\frac{3}{2}}=\left(\sqrt{9}\right)^3=27\)，并且 \(b^2=\left(\sqrt{3}\right)^2=3\). 因此
\(\left(\frac{b^2}{a}\right)^{\frac{1}{2}}=\left(\frac{3}{9}\right)^{\frac{1}{2}}=\left(\frac{1}{3}\right)^{\frac{1}{2}}=\frac{1}{\sqrt{3}}\)，
而 \(\left(\frac{1}{a}\right)^{-\frac{1}{2}}=\left(\frac{1}{9}\right)^{-\frac{1}{2}}=\frac{1}{\sqrt{1/9}}=\frac{1}{1/3}=3\).

所以原式为
\[
27\cdot\sqrt{3}\cdot\frac{1}{\sqrt{3}}\div 3=27\div 3=9
\]
故所求值为 \(9\).
\end{solution}

\begin{problem}
求函数 \(y=\sqrt{x-3}+\sqrt{5-x}\) 的自变量的取值范围.
\end{problem}

\begin{solution}
要使函数值为实数，两个平方根必须同时有意义，即两个被开方数都不小于零：\(x-3\geq 0\)，\(5-x\geq 0\). 这两个不等式分别给出 \(x\geq 3\) 和 \(x\leq 5\). 取它们的交集，得到 \(3\leq x\leq 5\)，且两个端点都使相应被开方数为零，因而可以取到. 所以自变量的取值范围是 \([3,5]\).
\end{solution}

\begin{problem}
求证：\(\cot\frac{\pi}{8}-\tan\frac{\pi}{8}=2\).
\end{problem}

\begin{solution}
令 \(\theta=\frac{\pi}{8}\). 因为 \(0<\theta<\frac{\pi}{2}\)，所以 \(\sin\theta\) 和 \(\cos\theta\) 均不为零，正切和余切都有定义. 由 \(\tan\theta=\frac{\sin\theta}{\cos\theta}\)、\(\cot\theta=\frac{\cos\theta}{\sin\theta}\)，有
\[
\cot\theta-\tan\theta
=\frac{\cos\theta}{\sin\theta}-\frac{\sin\theta}{\cos\theta}
=\frac{\cos^2\theta-\sin^2\theta}{\sin\theta\cos\theta}
=\frac{\cos 2\theta}{\frac{1}{2}\sin 2\theta}=2\cot 2\theta
\]
其中使用了 \(\cos^2\theta-\sin^2\theta=\cos 2\theta\) 和 \(2\sin\theta\cos\theta=\sin 2\theta\). 代入 \(\theta=\frac{\pi}{8}\)，得 \(2\theta=\frac{\pi}{4}\)，从而 \(\cot\frac{\pi}{8}-\tan\frac{\pi}{8}=2\cot\frac{\pi}{4}=2\times1=2\)，命题得证.
\end{solution}

\begin{problem}
要建造轮廓为矩形的猪舍四间，按实线部分砌墙，若墙的总长为 \(32\text{ 米}\)，问每间猪舍的长、宽为几米时面积最大？每间猪舍的最大面积是多少？

\FigureLayoutDeclare{img/1977/jiangsu_changzhou_city/q9_fig1.png}{6.2cm}{0bp 0bp 12bp 0bp}
\bitmapfigure[max width=6.2cm]{img/1977/jiangsu_changzhou_city/q9_fig1.png}
\end{problem}

\begin{solution}
设每间猪舍沿四间并列方向的长度为 \(x\text{ 米}\)，垂直于并列方向的宽度为 \(y\text{ 米}\). 图中实线的横向外墙由四个长度为 \(x\) 的部分组成，总长为 \(4x\)；纵向墙共有左、右两道外墙和三道间隔墙，共五道，每道长为 \(y\). 因而 \(4x+5y=32\)，并且 \(x>0\)、\(y>0\).

每间猪舍的面积为 \(S=xy\). 由墙长条件解得 \(y=\frac{32-4x}{5}\). 由于 \(y>0\)，还需有 \(0<x<8\). 代入面积式，得到
\[
S=x\cdot\frac{32-4x}{5}=-\frac{4}{5}x^2+\frac{32}{5}x
=-\frac{4}{5}(x-4)^2+\frac{64}{5}
\]
因为 \((x-4)^2\geq 0\)，所以 \(-\frac{4}{5}(x-4)^2\leq 0\)，从而 \(S\leq\frac{64}{5}\). 当且仅当 \(x=4\) 时等号成立，且此时 \(y=\frac{32-4\times4}{5}=\frac{16}{5}>0\)，满足可行条件. 因此 \(S_{\max}=\frac{64}{5}\).

所以每间猪舍的长为 \(4\text{ 米}\)，宽为 \(\frac{16}{5}\text{ 米}\)，每间猪舍的最大面积为 \(\frac{64}{5}\text{ 平方米}\).
\end{solution}

\begin{problem}
如图：割线 \(PO\) 交 \(\odot O\) 于 \(C\)，\(D\)，割线 \(PA\) 交 \(\odot O\) 于 \(A\)，\(B\)，

\begin{enumerate}
    \item 试证：\(PA\cdot PB=PC\cdot PD\)；
    \item 已知：\(\odot O\) 的半径为 \(1\)，\(PO=2\)，\(AB=1\)，求 \(PA\) 的长.
\end{enumerate}

\FigureLayoutDeclare{img/1977/jiangsu_changzhou_city/q10_fig1.png}{6.4cm}{10bp 32bp 2bp 4bp}
\bitmapfigure[max width=6.2cm]{img/1977/jiangsu_changzhou_city/q10_fig1.png}
\end{problem}

\begin{solution}
\begin{enumerate}
\item 连接 \(AC\) 和 \(BD\). 由图中点的次序，\(P\)，\(B\)，\(A\) 依次在同一直线上，\(P\)，\(D\)，\(O\)，\(C\) 依次在同一直线上，所以射线 \(PA\) 与 \(PB\) 同向，射线 \(PC\) 与 \(PD\) 同向，从而 \(\angle APC=\angle DPB\). 又因为射线 \(AP\) 与射线 \(AB\) 同向，所以 \(\angle PAC=\angle BAC\). 射线 \(DP\) 与射线 \(DC\) 反向，所以 \(\angle PDB=180^\circ-\angle CDB\).

四边形 \(ABDC\) 内接于圆，\(\angle BAC\) 和 \(\angle CDB\) 分别是从圆周上两点 \(A\)、\(D\) 看同一条弦 \(BC\) 所成的角；由于 \(A\)、\(D\) 是该圆内接四边形的对角顶点，二者互补，即 \(\angle BAC+\angle CDB=180^\circ\). 因此 \(\angle PAC=\angle BAC=180^\circ-\angle CDB=\angle PDB\). 由两个角分别相等，\(\triangle PAC\sim\triangle PDB\)，对应关系为 \(P\leftrightarrow P\)、\(A\leftrightarrow D\)、\(C\leftrightarrow B\). 所以
\(\frac{PA}{PD}=\frac{PC}{PB}\)，交叉相乘得 \(PA\cdot PB=PC\cdot PD\).

\item 因为 \(O\) 是圆心且半径为 \(1\)，所以 \(OC=OD=1\). 由图中点的次序 \(P\)，\(D\)，\(O\)，\(C\)，线段 \(PO\) 被近端交点 \(D\) 分成 \(PD\) 和 \(DO\)，所以 \(PD=PO-OD=2-1=1\)；而远端交点 \(C\) 在 \(O\) 的另一侧，所以 \(PC=PO+OC=2+1=3\).

设 \(PB=x\)，则 \(x>0\). 由图中点的次序 \(P\)，\(B\)，\(A\) 共线，且 \(AB=1\)，有 \(PA=PB+BA=x+1\). 代入（1）的割线乘积关系，并使用上面求得的 \(PC\)、\(PD\)，得到
\(x(x+1)=PA\cdot PB=PC\cdot PD=3\)，即 \(x^2+x-3=0\). 解得 \(x=\frac{-1\pm\sqrt{13}}{2}\). 其中 \(\frac{-1-\sqrt{13}}{2}<0\) 不符合线段长度 \(x>0\)，故 \(x=PB=\frac{\sqrt{13}-1}{2}\). 最后
\(PA=PB+AB=\frac{\sqrt{13}-1}{2}+1=\frac{\sqrt{13}+1}{2}\).
\end{enumerate}
\end{solution}

\begin{problem}
右图是某物体的二视图，按 \(1:1\) 的比例画出这个物体的直观图，并计算它的侧面积.
\FigureLayoutDeclare{img/1977/jiangsu_changzhou_city/q11_fig1.png}{6.0cm}{0bp 16bp 7bp 8bp}
\bitmapfigure[max width=6.0cm]{img/1977/jiangsu_changzhou_city/q11_fig1.png}
\end{problem}

\begin{solution}
俯视图是边长为 \(24\) 的正方形，且图中的中心与四个顶点相连. 这四条线表示四条侧棱在水平面上的投影，因此物体的顶点在正方形中心的正上方. 正视图是底边为 \(24\)、高为 \(16\) 的等腰三角形，说明棱锥的高为 \(16\). 所以该物体是底面边长为 \(24\)、高为 \(16\) 的正四棱锥.

按 \(1:1\) 比例作直观图时，先按边长 \(24\) 作出正方形底面并标出其中心；再从中心向上按比例取高 \(16\) 的顶点，最后把这个顶点分别与底面的四个顶点连接，即得到该正四棱锥的直观图.

设棱锥顶点为 \(S\)，底面中心为 \(O\)，任一底边的中点为 \(M\)，从 \(S\) 到该底边中点的斜高为 \(\ell\). 正四棱锥的高垂直于底面，所以 \(SO=16\)；又因为正方形中心到边的距离等于边长的一半，故 \(OM=\frac{24}{2}=12\). 在直角三角形 \(SOM\) 中，由勾股定理
\[
\ell=\sqrt{16^2+12^2}=\sqrt{256+144}=20
\]
每个侧面都是底为 \(24\)、高为 \(\ell=20\) 的等腰三角形，正四棱锥共有四个全等侧面，所以侧面积为
\[
4\times\frac{1}{2}\times 24\times 20=960
\]
故该物体的侧面积为 \(960\text{ 平方单位}\).
\end{solution}

\begin{problem}
已知：三点 \(A\left(-2,9\right)\)，\(B\left(4,-4\right)\)，\(C\left(10,8\right)\).

试求：

\begin{enumerate}
    \item 过 \(A\) 点且垂直于 \(BC\) 的直线方程；
    \item \(A\) 点到 \(BC\) 的距离.
\end{enumerate}
\end{problem}

\begin{solution}
\begin{enumerate}
\item 由 \(B(4,-4)\)、\(C(10,8)\) 得 \(BC\) 的斜率
\(k_{BC}=\frac{8-(-4)}{10-4}=\frac{12}{6}=2\). 所求直线与 \(BC\) 垂直，且 \(BC\) 不是水平线，所以所求直线的斜率为 \(k=-\frac{1}{k_{BC}}=-\frac{1}{2}\). 该直线经过 \(A(-2,9)\)，由点斜式
\(y-9=-\frac{1}{2}\left(x-(-2)\right)=-\frac{1}{2}(x+2)\). 两边乘以 \(2\)，得 \(2y-18=-x-2\)，整理得 \(x+2y-16=0\).

\item 由点 \(B(4,-4)\) 和斜率 \(2\)，直线 \(BC\) 的点斜式为 \(y-(-4)=2(x-4)\)，即 \(y+4=2x-8\)，整理得 \(2x-y-12=0\). 点 \((x_0,y_0)\) 到直线 \(u x+v y+w=0\) 的距离公式为
\(d=\frac{\lvert u x_0+v y_0+w\rvert}{\sqrt{u^2+v^2}}\). 本题中 \((x_0,y_0)=(-2,9)\)，\((u,v,w)=(2,-1,-12)\)，所以
\[
d=\frac{\left|2\times(-2)-9-12\right|}{\sqrt{2^2+(-1)^2}}=\frac{\lvert-25\rvert}{\sqrt{5}}=\frac{25}{\sqrt{5}}=5\sqrt{5}
\]
故 \(A\) 点到 \(BC\) 的距离为 \(5\sqrt{5}\).
\end{enumerate}
\end{solution}

\begin{problem}
如图：在直角坐标系中，\(OA_0=1\)，以 \(OA_0\) 为一条直角边画出第一个直角三角形 \(OA_0A_1\)，又以 \(OA_1\) 为一条直角边画出第二个直角三角形 \(OA_1A_2\)，再以 \(OA_2\) 为一条直角边画出第三个直角三角形 \(OA_2A_3\). 依次继续画下去，试求：

\begin{enumerate}
    \item \(A_0\) 点的极坐标；
    \item 第十二个直角三角形的面积（精确到 \(0.01\)）.
\end{enumerate}

已知 \(\lg 2=0.3010\)，\(\lg 3=0.4771\).

\begin{center}
反对数表

\begin{tabular}{c|cccccc}
\hline
\(m\) & 0 & 1 & 2 & 3 & 4 & 5 \\
\hline
\(80\) & \(0.6310\) & \(0.6324\) & \(0.6339\) & \(0.6353\) & \(0.6368\) & \(0.6383\) \\
\(81\) & \(0.6457\) & \(0.6471\) & \(0.6486\) & \(0.6501\) & \(0.6516\) & \(0.6531\) \\
\(82\) & \(0.6607\) & \(0.6622\) & \(0.6637\) & \(0.6653\) & \(0.6668\) & \(0.6683\) \\
\(83\) & \(0.6761\) & \(0.6776\) & \(0.6792\) & \(0.6808\) & \(0.6823\) & \(0.6839\) \\
\(84\) & \(0.6918\) & \(0.6934\) & \(0.6950\) & \(0.6966\) & \(0.6982\) & \(0.6998\) \\
\hline
\end{tabular}
\end{center}

\begin{center}
\bitmapfigure[float=inline,max width=0.86\linewidth]{img/1977/jiangsu_changzhou_city/q13_fig1.png}
\end{center}
\end{problem}

\begin{solution}
在极坐标 \((r,\theta)\) 中，\(r\) 表示点到极点的距离，\(\theta\) 表示从极轴正半轴到该点所在射线的转角. 由图可知，点 \(A_0\) 在 \(x\) 轴正半轴上，且 \(OA_0=1\)，所以其极径为 \(1\)、极角为 \(0\)，即 \(A_0\) 的极坐标为 \(\left(1,0\right)\).

由图中各射线的标注，相邻两条射线 \(OA_{n-1}\)、\(OA_n\) 的夹角均为 \(\frac{\pi}{6}\). 又因为以 \(OA_{n-1}\) 为一条直角边作三角形，直角在 \(A_{n-1}\) 处，所以在直角三角形 \(OA_{n-1}A_n\) 中，\(OA_n\) 是斜边，\(OA_{n-1}\) 是角 \(\frac{\pi}{6}\) 的邻直角边. 设 \(r_n=OA_n\)，则
\[
\cos\frac{\pi}{6}=\frac{OA_{n-1}}{OA_n}=\frac{r_{n-1}}{r_n}
\]
由于 \(\cos\frac{\pi}{6}=\frac{\sqrt{3}}{2}\)，所以
\[
r_n=\frac{r_{n-1}}{\cos\frac{\pi}{6}}=\frac{2}{\sqrt{3}}r_{n-1}
\]
初始值为 \(r_0=OA_0=1\)，反复使用上式可得
\[
r_{11}=\left(\frac{2}{\sqrt{3}}\right)^{11}
\]
第十二个三角形是 \(OA_{11}A_{12}\). 它在 \(A_{11}\) 处为直角，且 \(OA_{11}=r_{11}\). 在顶点 \(O\) 处，\(OA_{11}\) 与 \(OA_{12}\) 的夹角为 \(\frac{\pi}{6}\)，所以以 \(OA_{11}\) 为邻直角边、以 \(A_{11}A_{12}\) 为对边，有
\[
\tan\frac{\pi}{6}=\frac{A_{11}A_{12}}{OA_{11}}=\frac{A_{11}A_{12}}{r_{11}}
\]
从而
\[
A_{11}A_{12}=r_{11}\tan\frac{\pi}{6}=\frac{r_{11}}{\sqrt{3}}
\]
因此第十二个三角形的面积为
\[
S_{12}=\frac{1}{2}r_{11}\cdot\frac{r_{11}}{\sqrt{3}}
=\frac{1}{2\sqrt{3}}\left(\frac{2}{\sqrt{3}}\right)^{22}
=\frac{2^{21}}{3^{23/2}}
\]
利用题给对数，
\[
\lg S_{12}=21\lg 2-\frac{23}{2}\lg 3
=21\times0.3010-\frac{23}{2}\times0.4771=0.83435
\]
查反对数表，尾数 \(0.8340\)、\(0.8350\) 位于第 \(83\) 行的第 \(4\)、\(5\) 列，表中数值分别为 \(0.6823\)、\(0.6839\). 由于特征数为 \(0\)，对应的反对数分别为 \(6.823\)、\(6.839\). 又因为 \(0.83435=0.8340+0.35\times0.0010\)，在线性插值下
\[
S_{12}\approx 6.823+0.35\left(6.839-6.823\right)=6.8286
\]
按题意精确到 \(0.01\)，得
\[
S_{12}\approx 6.83
\]
故第十二个直角三角形的面积约为 \(6.83\text{ 平方单位}\).
\end{solution}
